\section{Introduction}\label{sec:intro}

This manuscript reports the first results of the \emph{Loewner} investigation
of the SUSY-positivity program. It is written to be read on its own, and its
introduction is deliberately slow: the results in Sections~\ref{sec:decomposition}
to~\ref{sec:contraction} are short and elementary once the objects are in
view, and the work of the investigation was largely in getting them into view.
The chain of logic runs from Weil's criterion, through the shifted family and
its reading as a filter, through a proposal to realize that filter by a
conformal evolution, to a factorization of the filter that says what any
realization would have to produce --- and then to a computation that confirms
the factorization is what it claims to be.

\subsection{Weil's criterion, and where the arithmetic hides}

Riemann's zeta function has a completed form
$\xi(s)=\frac12s(s-1)\pi^{-s/2}\Gamma(s/2)\zeta(s)$, entire, real on the real
axis, symmetric under $s\mapsto1-s$, whose zeros are the nontrivial zeros
$\rho$ of $\zeta$. The Riemann hypothesis places them all on the critical line
$\re s=\frac12$. Weil's criterion restates this as the positivity of a
quadratic form. Take a test function $f$ supported in a window
$I_L=(-\frac L2,\frac L2)$ of the real line, extend it by zero to $F$, and
write $\gamma_\rho=-i(\rho-\frac12)$, so that $\gamma_\rho$ is real exactly when
$\rho$ is on the line. The \emph{Weil form} is
\begin{equation}\label{eq:intro-spectral}
 Q_L[f]=\sum_\rho\widehat F(\gamma_\rho)\,\overline{\widehat F(\overline{\gamma_\rho})},
\end{equation}
and the criterion says that RH holds if and only if $Q_L[f]\geq0$ for every
$L$ and every $f$. Under RH each term is $|\widehat F(\gamma_\rho)|^2$ and the
positivity is obvious; a zero off the line pairs with its reflection to produce
a term that can be made negative by a suitable $f$.

The reason this is not a tautology is the explicit formula. The same quantity
$Q_L[f]$ has a second expression in which the zeros do not appear at all: an
\emph{archimedean} term built from the digamma function, a \emph{pole} term
from the factor $s(s-1)$, and a \emph{prime} term,
$-2\sum_{m\log p<L}(\log p)\,p^{-m/2}\,\re\ip F{U_{m\log p}F}$, in which $F$ is
correlated with its own translates by the logarithms of prime powers. The
prime-side expression is what one can compute, and what a construction could
hope to produce; the zero-side expression is what is positive. The whole
difficulty of the criterion is that positivity is manifest on one side and
invisible on the other, and the two sides are the same number. The cancellation
between them is extreme: at $L=\log 3$ the smallest value of $Q_L$ on
normalized $f$ is about $5\times10^{-8}$, at $L=\log 7$ about
$10^{-27}$, while the individual terms are of order one. Every investigation of
this program is, in one way or another, a search for a presentation of $Q_L$ in
which positivity is structural and the arithmetic is nevertheless present.

\subsection{The shifted family: RH as the stability of a filter}

The Wilson-lines investigation \cite{WilsonLines} studies a one-parameter
deformation of the criterion. Write $s=\frac12+p$ and, for $0<\omega\leq\frac12$,
\begin{equation}\label{eq:intro-transfer}
 K_\omega(p)=\frac{\xi(\tfrac12+p-\omega)}{\xi(\tfrac12+p+\omega)} .
\end{equation}
Three facts organize everything that follows, and all three are elementary.
First, $K_\omega$ is \emph{unimodular on the imaginary axis}: the functional
equation and the reality of $\xi$ make numerator and denominator complex
conjugates at $p=i\tau$, so $|K_\omega(i\tau)|=1$ for every real $\tau$. In the
language of signal processing $K_\omega$ is an \emph{all-pass filter}: it
changes the phase of every frequency and the amplitude of none. Second,
$K_\omega$ is the symbol of a \emph{causal} operator: realized as a Laplace
transform on a line $\re p=\eta$ to the right of all its singularities, it is
$\int_0^\infty e^{-pt}k_\omega(t)\,dt$ for a kernel $k_\omega$ supported on
$t\geq0$, and the operator $V_\omega f=k_\omega*f$ never moves mass to the left.
Third, the derivative at $\omega=0$ is the Weil form. Compress $V_\omega$ to the
window, $V_{\omega,L}=P_LV_\omega P_L$ with $P_L$ the restriction to $I_L$, and
form the \emph{contraction defect} $D_{\omega,L}=I-V_{\omega,L}^*V_{\omega,L}$,
which measures how far the compression is from an isometry. Then
\begin{equation}\label{eq:intro-firstorder}
 \norm f^2-\norm{V_{\omega,L}f}^2=\ip f{D_{\omega,L}f}=2\omega\,Q_L[f]+O(\omega^2).
\end{equation}
So Weil positivity on $I_L$ is the statement that, to first order in the shift,
the compressed filter contracts.

The three facts combine into a dichotomy \cite[Proposition~7.2]{WilsonLines}
that is the hinge of the present manuscript. Fix $\omega$. On the whole line an
all-pass filter is unitary, and a compression of a unitary is a contraction; but
$V_\omega$ is unitary on $L^2(\R)$ only if the causal realization can be moved
onto the imaginary axis, which requires $K_\omega$ to have no poles in the right
half-plane. The poles of $K_\omega$ are at $p=\rho-\frac12-\omega$, so this is
the statement that no zero has $\re\rho>\frac12+\omega$. If there is such a
zero, $k_\omega$ contains a growing exponential and the compressions are
unbounded as $L$ grows. Hence
\begin{equation}\label{eq:intro-dichotomy}
 \sup_L\norm{V_{\omega,L}}\in\{1,\infty\},\qquad
 \sup_L\norm{V_{\omega,L}}=1\iff\text{no zero of $\xi$ lies further than $\omega$ from the line.}
\end{equation}
The shifted family is a graded criterion, one zero-free strip per shift, and
the Riemann hypothesis is the statement that, for every $\omega>0$, the causal
filter $k_\omega$ is \emph{passive}: every window compression is a contraction.
Under RH the defect has a physical meaning, the flux identity
\cite[Proposition~4.5]{WilsonLines}: $\ip f{D_{\omega,L}f}$ is exactly the mass
of $V_\omega f$ that has been pushed past the leading end of the window. The
form $Q_L$ is the first-order version of that flux.

\subsection{Realizations, and the two proposals}

The bulk-realization program asks for a system --- a field theory, a
statistical model, a conformal evolution --- in which the transfer, or the
form, appears as an observable whose positivity is structural, so that the
arithmetic would be carried by the system and its positivity by a general
theorem (unitarity, reflection positivity, subordination). The dichotomy sets
the standard any such realization must meet: it must be exactly critical,
unitary on the whole line and contracting on every window by a margin
$\omega m_L$ that collapses superexponentially in $L$ \cite[Section~8]{WilsonLines}.

The Wilson-lines proposal \cite{WilsonLines} took the transfer to be the
deformation flow of an open Wilson line, because the shifted family obeys a
flow equation $\partial_\omega V_{\omega,L}=-G_{\omega,L}V_{\omega,L}$ of the
shape produced by varying a line's contour, with the non-Abelian Stokes theorem
supplying the coefficient. Two of that manuscript's findings bear directly on
what follows. The shift direction is \emph{abelian}: compression to a window
is multiplicative on causal kernels, so the transfers and their generators lie
in one commutative algebra --- the algebra of translation-invariant causal
operators on the line, lower-triangular Toeplitz matrices in any discretization
--- and what is abelian is precisely translation invariance along the
arithmetic coordinate \cite[Theorem~5.2]{WilsonLines}. And the
two-parameter flow in the shift--length plane is flat
\cite[Proposition~9.4]{WilsonLines}. The inherited note \cite{LoewnerNote}
from which this investigation was opened sharpened both statements, in response
to the question of whether these were obstructions or selections. They are
selections. Flatness is a tautology for any single-valued family
$V(\alpha)$ with $dV=\mathcal AV$ --- the curvature of $\mathcal A=dV\,V^{-1}$
vanishes identically --- so it constrains nothing about a realization; in the
non-Abelian Stokes variation the field strength is the \emph{coefficient}
$\mathcal A$, not the curvature, and the object to match is the generator
$G_{\omega,L}$ itself. Commutativity, likewise, is a symmetry requirement on the
insertions that implement the shift (they must act as functions of one
operator, uniformly along the line), not an impossibility. Along a general path
in the $(\omega,L)$ plane the transport is a genuinely ordered product, because
the endpoint generator $B_L$ and the bulk generator $G_{\omega,L}$ generate a
triangular, not commutative, algebra.

The Loewner proposal, made at that point, was to construct the paths
independently: to take the arithmetic coordinate $x=\log r$ as the time of a
whole-plane Loewner evolution, whose hull grows from $0$ toward $\infty$ with
logarithmic capacity $t$, and to let a simple boundary driving function
produce the ordered, infinite-dimensional, structurally contractive family that
the shift direction lacks. The affinities are real: Loewner time is
log-conformal-radius, which is the arithmetic coordinate exactly; the marked
points $0,\infty$ have the dilations and inversions of the Weil form's exact
symmetry group as their stabilizer; the Loewner equation is a tip equation
with one scalar coefficient, like the endpoint variation of the transfer; and
composition with univalent self-maps is contractive on Hardy-type spaces by
Littlewood subordination, with a positive defect and no arithmetic in it.

\subsection{What a realization would have to produce: the factorization}

To test either proposal one must know the kernel $k_\omega$ concretely, and it
turns out to have a structure that neither proposal anticipated. Factor the
completed zeta function as $\xi(u)=\frac12u(u-1)\Lambda(u)$ with
$\Lambda(u)=\pi^{-u/2}\Gamma(u/2)\zeta(u)$. The transfer splits
correspondingly into a ratio of $\Lambda$'s and a ratio of quadratics,
$K_\omega=R_\omega\Kt_\omega$, and the two factors have opposite characters.

The ratio $\Kt_\omega(p)=\Lambda(s-\omega)/\Lambda(s+\omega)$ is
\emph{completely monotone} on $p>\frac12+\omega$: all its derivatives alternate
in sign. By Bernstein's theorem a completely monotone function is the Laplace
transform of a positive measure, so the causal kernel $\kt_\omega$ of
$\Kt_\omega$ is \emph{nonnegative}. Moreover it is explicit. The $\Gamma$-ratio
$\pi^\omega\Gamma(\frac{s-\omega}2)/\Gamma(\frac{s+\omega}2)$ is a Beta integral
whose kernel is
$k^\Gamma_\omega(x)=\frac{2\pi^\omega}{\Gamma(\omega)}(2\sinh x)^{\omega-1}e^{x/2}$,
a continuous delay distributed as $x=-\frac12\log U$ with
$U\sim\mathrm{Beta}(\frac14-\frac\omega2,\omega)$; and the $\zeta$-ratio
$\zeta(s-\omega)/\zeta(s+\omega)$ is a Dirichlet series whose coefficients
$\ct_n=n^{\omega-1/2}\prod_{p\mid n}(1-p^{-2\omega})$ are positive because the
Euler product makes each local factor a geometric series with positive terms.
Its kernel is a comb of atoms at the integer logarithms, and
$\kt_\omega=k^\Gamma_\omega*\sum_n\ct_n\delta_{\log n}$: a sum of Beta-shaped
spikes, one at every $\log n$, with positive weights. This is the
\emph{Markov part} of the transfer, a positive kernel, though of infinite
mass --- it grows like $e^{(\frac12+\omega)x}$, because $\Lambda(s-\omega)$ has
the pole of $\zeta$ at $s-\omega=1$, and that pole is the prime number theorem's
main term.

The rational factor $R_\omega$ is a ratio of two Blaschke factors, unimodular
on the axis, and it is \emph{forced}: given the Markov part, it is the unique
minimal rational all-pass function that cancels the poles and zeros of
$\Kt_\omega$ at the poles of $\Lambda$, so its coefficients carry nothing but
the positions $0,1$ of those poles and the shift. Absorbing the harmless half
of it into the Markov part gives the form we use throughout:
\begin{equation}\label{eq:intro-decomp}
 K_\omega=B_b\,\Kh_\omega,\qquad B_b(p)=\frac{p-b}{p+b},\qquad b=\tfrac12+\omega,
 \qquad \Kh_\omega\ \text{completely monotone on }p>b,
\end{equation}
and at the level of kernels
\begin{equation}\label{eq:intro-ema}
 k_\omega=\kh_\omega-2\,\EMA_b[\kh_\omega],\qquad
 \EMA_b[g](x)=b\int_0^xe^{-b(x-y)}g(y)\dd y,\qquad \kh_\omega\geq0 .
\end{equation}
\emph{The impulse response of the shifted Weil transfer is a positive measure
minus twice its exponential moving average}, at the rate set by the pole of
$\zeta$. The subtraction is exact and unconditional: it maps the main term
$Ce^{bx}$ of the positive part to $Ce^{-bx}$, so the growth is cancelled
identically, and what remains of $k_\omega$ is the fluctuation of $\kh_\omega$
about its main term --- which is to say, the zeros. Under RH the remainder
decays like $e^{-\omega x}$; a zero at distance $\delta>\omega$ from the line
leaves a growing $e^{(\delta-\omega)x}$. This is the dichotomy
\eqref{eq:intro-dichotomy} at the level of the kernel, and it says where the
arithmetic hides: not in the small rational correction, whose residues are
fixed by the functional equation, but in the fine structure of a positive
measure whose Laplace transform has poles at the zeros behind its abscissa of
convergence.

\subsection{What the pieces are, and what can make them}

Each piece of the positive part has a name in probability or geometry, and the
names decide the two proposals. The Beta law is \emph{Bessel additivity}: with
$a=\frac12-\omega$ and $b=\frac12+\omega$, $U=Z_a/(Z_a+Z_{2\omega})$ for
independent squared Bessel processes of dimensions $a$ and $2\omega$, whose sum
is a squared Bessel process of dimension $b$. The two exponents of the rational
factor are the two Bessel dimensions, and the delay test's archimedean quarter
is half a Bessel dimension. This is a statement about Bessel processes, which
drive radial Loewner chains but are not themselves chains; whether a radial
chain has this law as a squared-radius ratio is the one question in the Loewner
direction that remains open.

The comb cannot come from a Loewner chain in any reading of the variables. If
the arithmetic coordinate is Loewner time, the only translation-invariant
filter a rotation-invariant chain induces on analytic observables is the pure
delay $e^{-tp}$ --- a rotation average of a composition operator is a dilation.
If the arithmetic coordinate is the radial coordinate $\log|z|$, the chain's
transport of exterior points is a pure delay in the far field,
$\log|g_t(z)|=\log|z|-t+O(|z|^{-1})$, and translation invariance forces the
induced kernel to equal its far-field limit. Echoes at the fixed separations
$\log n$ with position-independent weights are available in neither reading.
What the comb \emph{is} is a semigroup structure: the operator
$C_\omega=\sum_n\ct_n\mu_n$ with $\mu_nf(x)=f(x-\log n)$ and
$\mu_n\mu_m=\mu_{nm}$, a Dirichlet series in the multiplicative semigroup
$\N^\times$ acting on the ray by dilations --- the Hecke operator of the
Bost--Connes system, with coefficients that are, to first order in $\omega$,
$2\omega\Lambda(n)n^{-1/2}$: the von Mangoldt function.

And at the endpoint $\omega=\frac12$ the whole positive part is a classical
geometric object: $\Kh_{1/2}=\Lambda(p)/\Lambda(p+1)$ is the Eisenstein
scattering matrix of the modular surface $PSL(2,\mathbb Z)\backslash\mathbb H$,
the comb weights are Euler's totient $\varphi(n)/n$ --- the coset count in the
cusp's Fourier expansion that produces $\zeta(2s-1)/\zeta(2s)$ in the first
place --- and $K_{1/2}=B_1\Kh_{1/2}$ is that scattering matrix with its pole at
$s=1$ removed by one Blaschke factor, which is exactly the Lax--Phillips
modification. So the Markov part has, at one shift, a realization in which the
comb is a coset sum and the correction is the removal of the residual
spectrum. For $\omega<\frac12$ no Eisenstein construction produces the offset
$2\omega$, and the interior of the family has no such realization from that
source.

\subsection{What the decomposition is not, and what it is}

It is not a positivity mechanism, and the dichotomy says why with no
computation. The positive part $\Kh_\omega$ has a causal operator $\widehat
V_\omega$ of norm greater than one on every window (its kernel grows), the
Blaschke factor is a contraction, and their product is the transfer; the
contraction of $V_{\omega,L}$ for all $L$ is a zero-free strip by
\eqref{eq:intro-dichotomy}, so no structural property of $\kh_\omega$ short of
the location of the poles can imply it. Any passivity theorem for filters of
the form ``positive measure minus twice its EMA'' would have to see the poles.
We record this plainly because it is the honest end of one line of thought:
positivity of the Markov part is real, explicit and cheap, and it proves
nothing about the zeros.

What the decomposition is, is three things. It is a \emph{computational tool}:
on a window of length $L$ only the atoms $\log n<L$ contribute, so the
transfer, its compression, its norm and its defect are assembled exactly from
$(2\sinh x)^{\omega-1}e^{x/2}$, the numbers $\ct_n$ for $n<e^L$, and two
exponentials --- no $\xi$, no $\zeta$ on a line, no zeros --- in precisely the
way the Weil form is assembled from the explicit formula. It is a
\emph{dictionary}: Bessel additivity for the archimedean factor, the Hecke
semigroup for the comb, one Blaschke factor at the pole of $\zeta$ for the
correction, the modular surface at the endpoint. And it is a statement of the
program's criticality in filter language: the shifted Weil transfer is a
lossless filter whose impulse response is a positive measure minus twice its
exponential moving average, and the Riemann hypothesis is the statement that
this filter is stable at every $\omega>0$.

\subsection{The first computation}

The tool was then built and tested. Section~\ref{sec:contraction} assembles
$V_{\omega,L}$ at the first horizon $L=\log3$ in an orthonormal sine basis,
from the kernel alone, and measures its contraction data. The compression is a
strict contraction at every shift computed, and the smallest eigenvalue of its
defect satisfies $\lambda_{\min}(D_{\omega,L})/2\omega\to m_L$, the margin of
the Weil form in the same basis, with agreement to $0.1\%$ at $\omega=0.01$:
the transfer side and the form side of the program, assembled from different
presentations of the explicit formula, agree at first order. The deviations at
larger $\omega$ are then accounted for exactly. On the exact compression the
defect expands as $D_{\omega,L}=2\omega Q_L-\omega^2(2Q_L^2+[Q_L,A_L])+O(\omega^3)$,
where $A_L$ is the skew part of the compressed generator, and on the minimizer
of $Q_L$ the commutator term vanishes, so the only first-order correction to
\eqref{eq:intro-firstorder} at the minimizer is $-\omega m_L^2$. The positive
excess seen in a finite basis is a truncation term, the energy of $A_Lf$
outside the basis, and it dies under nesting exactly as predicted. Finally the
mass of the assembled kernel beyond the horizon has, to first order in
$\omega$, a closed form from the logarithmic derivative of $\xi$ alone ---
poles, digamma and primes, no zeros --- and the assembly matches it to five
digits.

\subsection{Organization and status}

Section~\ref{sec:conventions} fixes conventions and quotes, with proofs where
they are short, the facts about the shifted family that the rest uses.
Section~\ref{sec:decomposition} proves the factorization and its reading.
Section~\ref{sec:realizations} identifies the pieces and proves the negative
Loewner results. Section~\ref{sec:contraction} presents the assembly, the
second-order theory of the defect, the first-order tail, and the numerical
results. Section~\ref{sec:outlook} ranks what remains. Appendix~\ref{app:numerics}
documents the programmes and records and the numerical lessons;
Appendix~\ref{app:provenance} records provenance and the use of language
models. Every statement is labelled as a written proof, a labelled numerical
computation, or a reading; nothing assumes the Riemann hypothesis except where
a statement says so.
